\section{Examples}

To illustrate the ideas of the API, we present several short examples that demonstrate how the API can be used. The
full details of the interface are presented in the next section. To prevent confusion, we refer to the application
process or binary that is being modified as the mutatee, and the program that uses the API to modify the application as
the mutator. The mutator is a separate process from the application process.

The examples in this section are simple code snippets, not complete programs. Appendix A - Complete Examplesprovides
several examples of complete Dyninst programs.}

\subsection{Instrumenting a function}

A mutator program must create a single instance of the class BPatch. This object is used to access functions and
information that are global to the library. It must not be destroyed until the mutator has completely finished using
the library. For this example, we assume that the mutator program has declared a global variable called bpatch of
class BPatch.

All instrumentation is done with a BPatch_addressSpace object, which allows us to write codes that work for both
dynamic and static instrumentation. During initialization we use either BPatch_process to attach to or create a
process, or BPatch_binaryEdit to open a file on disk. When instrumentation is completed, we will either run the
BPatch_process, or write the BPatch_binaryEdit back onto the disk.

The mutator first needs to identify the application to be modified. If the process is already in execution, this can
be done by specifying the executable file name and process id of the application as arguments in order to create an
instance of a process object:

\code{BPatch_process *appProc = bpatch.processAttach(name, processId);}

This creates a new instance of the \code{BPatch_process} class that refers to the existing process. It had no effect on the
state of the process (i.e., running or stopped). If the process has not been started, the mutator specifies the
pathname and argument list of a program it seeks to execute:

\code{BPatch_process *appProc = bpatch.processCreate(pathname, argv);}

If the mutator is opening a file for static binary rewriting, it executes:

\code{BPatch_binaryEdit *appBin = bpatch.openBinary(pathname);}

The above statements create either a \code{BPatch_process} object or \code{BPatch_binaryEdit} object, depending on whether Dyninst
is doing dynamic or static instrumentation. The instrumentation and analysis code can be made agnostic towards static
or dynamic modes by using a BPatch_addressSpace object. Both \code{BPatch_process} and \code{BPatch_binaryEdit} inherit from
\code{BPatch_addressSpace}, so we can use cast operations to move between the two:

\code{BPatch_process *appProc = static_cast<BPatch_process *>(appAddrSpace);}

or

\code{BPatch_binaryEdit *appBin = static_cast<BPatch_binaryEdit *>(appAddrSpace);}

Similarly, all instrumentation commands can be performed on a \code{BPatch_addressSpace} object, allowing similar codes to be
used between dynamic instrumentation and binary rewriting:

\code{BPatch_addressSpace *app = appProc;}

or

\code{BPatch_addressSpace *app = appBin;}

Once the address space has been created, the mutator defines the snippet of code to be inserted and identifies where the
points should be inserted.

If the mutator wants to instrument the entry point of \code{InterestingProcedure}, it should get a
\code{BPatch_function} from the application\'s \code{BPatch_image}, and get the entry 
\code{BPatch_point} from that function:

\code{
  std::vector<BPatch_function *> functions;
  std::vector<BPatch_point *> *points;

  BPatch_image *appImage = app->getImage();
  appImage->findFunction(“InterestingProcedure”, functions);
  points = functions[0]->findPoint(BPatch_locEntry);
}

The mutator also needs to construct the instrumentation that it will insert at the \code{BPatch_point}.
It can do this by allocating an integer in the application to store instrumentation results, and
then creating a \code{BPatch_snippet} to increment that integer:

\code{
  BPatch_variableExpr *intCounter = app->malloc(*(appImage->findType("int")));
  BPatch_arithExpr addOne(BPatch_assign, *intCounter,}
  BPatch_arithExpr(BPatch_plus, *intCounter, BPatch_constExpr(1)));
}

The mutator can set the \code{BPatch_snippet} to be run at the \code{BPatch_point} by
executing an \code{insert­Snippet} call:

\code{app->insertSnippet(addOne, *points);}

Finally, the mutator should either continue the mutate process and wait for it to finish,
or write the resulting binary onto the disk, depending on whether it is doing dynamic
or static instrumentation:

\code{
  appProc->continueExecution();

  while (!appProc->isTerminated()) {
    bpatch.waitForStatusChange();
  }
}

or

\code{appBin->writeFile(newPath);}

A complete example can be found in Appendix A - Complete Examples.

\subsection{Binary Analysis}

This example will illustrate how to use Dyninst to iterate over a function\'s control flow graph and
inspect instructions. These are steps that would usually be part of a larger data flow or control flow analysis.
Specifically, this example will collect every basic block in a function, iterate over them, and count the number of
instructions that access memory.

Unlike the previous instrumentation example, this example will analyze a binary file on disk. Bear in mind, these
techniques can also be applied when working with processes. This example makes use of InstructionAPI, details of
which can be found in the InstructionAPI Reference Manual.

Similar to the above example, the mutator will start by creating a \code{BPatch_object} and opening a file to
operate on:

\code{
  BPatch bpatch;
  BPatch_binaryEdit *binedit = bpatch.openFile(pathname);
}

The mutator needs to get a handle to a function to do analysis on. This example will look up a function by name;
alternatively, it could have iterated over every function in \code{BPatch_image} or \code{BPatch_module}:

\code{
  BPatch_image *appImage = binedit->getImage();
  std::vector<BPatch_function *> funcs;
  image->findFunction(“InterestingProcedure”, funcs);
}

A function\'s control flow graph is represented by the \code{BPatch_flowGraph} class. The
\code{BPatch_flowGraph} contains, among other things, a set of \code{BPatch_basicBlock} objects
connected by \code{BPatch_edge} objects. This example will simply collect a list of the basic blocks
in \code{BPatch_flowGraph} and iterate over each one:

\code{
  BPatch_flowGraph *fg = funcs[0]->getCFG();
  std::set<BPatch_basicBlock *> blocks;
  fg->getAllBasicBlocks(blocks);

  for(BPatch_basicBlock *b : blocks) {
    std::vector<Dyninst::InstructionAPI::Instruction> insns;
    block->getInstructions(insns);
  }
}

Given an \code{Instruction} object, which is described in the InstructionAPI Reference Manual, we can
query for properties of this instruction. InstructionAPI has numerous methods for inspecting the
memory accesses, registers, and other properties of an instruction. This example simply checks whether
this instruction accesses memory:

\code{
  int insns_access_memory = 0;
  for(auto const& insn : insns) {
    if(insn.readsMemory() || insn->writesMemory()) {
      insns_access_memory++;
    }
  }
}

\subsection{Instrumenting Memory Accesses}

There are two snippets useful for memory access instrumentation: \code{BPatch_effectiveAddressExpr} and
\code{BPatch_bytesAccessedExpr}. Both have nullary constructors; the result of the snippet depends on the
instrumentation point where the snippet is inserted. \code{BPatch_effectiveAddressExpr} has type
\code{void*}, while \code{BPatch_bytesAccessedExpr} has type \code{int}.

These snippets may be used to instrument a given instrumentation point if and only if the point has memory access
information attached to it. In this release the only way to create instrumentation points that have memory access
information attached is via \code{BPatch_function.findPoint(const std::set<BPatch_opCode>&)}. For
example, to instrument all the loads and stores in a function named \code{InterestingProcedure} with a call to
\code{printf}, one may write:

\code{
  BPatch_addressSpace *app = ...;
  BPatch_image *appImage = proc->getImage();

  // We're interested in loads and stores
  std::set<BPatch_opCode> axs;
  axs.insert(BPatch_opLoad);
  axs.insert(BPatch_opStore);

  // Scan the function InterestingProcedure and create instrumentation points
  std::vector<BPatch_function*> funcs;
  appImage->findFunction(“InterestingProcedure”, funcs);

  std::vector<BPatch_point*>* points = funcs[0]->findPoint(axs);
  
  // Create the printf function call snippet
  std::vector<BPatch_snippet*> printfArgs = {
    new BPatch_constExpr("Access at: \%p.\n"),
    new BPatch_effectiveAddressExpr()
  };

  // Find the printf function
  std::vector<BPatch_function *> printfFuncs;
  appImage->findFunction("printf", printfFuncs);

  // Construct the function call snippet
  BPatch_funcCallExpr printfCall(*(printfFuncs[0]), printfArgs);

  // Insert the snippet at the instrumentation points
  app->insertSnippet(printfCall, *points);
}
